% BEGIN PROFESSOR FEEDBACK PF19: Results - RQ1
\section{RQ1: Incremental Signal of Temporal Motif Fingerprints}

Adding local motif-role features increases average precision in all six
HuggingFace comparisons. Three of these increases are statistically significant.
Figure~\ref{fig:res_rq1_ablation} compares the baseline, motif-only, and combined
feature sets in the 40-role analysis. Table~\ref{tab:res_rq1} shows the same
comparison for all three samples. The addresses and data splits
remain identical within each comparison, so only the feature set changes.

\begin{figure}[ht]
  \centering
  \ResGraphic{0.48}{rq1_ablation_controlled_rep3}
  \ResGraphic{0.48}{rq1_ablation_prevalence_rep3}
  \caption{RQ1 feature comparison in the 40-role HuggingFace analysis using the
  default timestamp ordering. The controlled cohort is shown on the left and
  the prevalence cohort on the right.}
  \label{fig:res_rq1_ablation}
\end{figure}

\IfFileExists{generated/results_rq1.tex}{\begin{table}[ht]
\centering
\caption{RQ1: average precision of each feature view at the selected window (canonical ordering), the combined-minus-baseline difference with its 95\,\% confidence interval, the permutation $p$-value, and the decision.}
\label{tab:res_rq1}
\footnotesize
\setlength{\tabcolsep}{4pt}
\begin{tabular}{@{}l l r r r r c c l@{}}
\toprule
\textbf{Cohort} & \textbf{$\delta$} & \textbf{Base} & \textbf{Motifs} & \textbf{Comb.} & \textbf{$\Delta$AP} & \textbf{95\,\% CI} & \textbf{$p$} & \textbf{Decision} \\
\midrule
Controlled & 5 min & 0.603 & 0.471 & 0.533 & -0.071 & [-0.219, 0.033] & 0.663 & not rejected \\
Prevalence & 5 min & 0.682 & 0.449 & 0.660 & -0.022 & [-0.109, 0.102] & 0.634 & not rejected \\
\bottomrule
\end{tabular}
\end{table}
}{\ResPendingTable{RQ1 feature-view comparison.}{tab:res_rq1}}

Three of the six improvements are statistically significant. Both cohorts are
significant in one 32-role sample, while only the controlled cohort is
significant in the other. Neither cohort is significant in the 40-role
analysis. In that analysis, average precision rises from 0.7351 to
0.7587 in the controlled cohort. The difference is 0.0237, with a 95\,\%
confidence interval of $[-0.0142,0.0615]$ and $p=0.1330$. In the prevalence
cohort, it rises from 0.5986 to 0.6146. Here the difference is 0.0161, with a
confidence interval of $[-0.0342,0.0653]$ and $p=0.2659$. The direction remains
positive under all seven timestamp orderings, but the confidence intervals
include zero.

The motif-only model has lower AP than the non-motif baseline in every
sample. The
role-by-role fingerprint performs better than using only the total number of
motifs. In the 40-role analysis, its AP advantage over motif volume is 0.1173 in
the controlled cohort and 0.1528 in the prevalence cohort. Both differences
remain significant after correction for multiple comparisons. Once the baseline
is supplemented with motif composition and volume, adding the full role-by-role
fingerprint does not give a significant improvement.

\IfFileExists{generated/results_rq1_representations.tex}{\begin{table}[ht]
\centering
\caption{Secondary RQ1 representation comparisons at the development-selected window using the default timestamp ordering. The first comparison tests fingerprint shape beyond total motif count. The second tests the combined model against a baseline augmented with motif composition and volume. P-values use Benjamini--Hochberg correction within the comparison family for each sample.}
\label{tab:res_rq1_representations}
\scriptsize
\setlength{\tabcolsep}{5pt}
\begin{tabular}{@{}r l p{3.4cm} r r r c r l@{}}
\toprule
\textbf{Sample} & \textbf{Cohort} & \textbf{Comparison} & \textbf{Ref.} & \textbf{Alt.} & \textbf{$\Delta$AP} & \textbf{95\,\% CI} & \textbf{$p_{\mathrm{BH}}$} & \textbf{Significant?} \\
\midrule
1 & Controlled & combined vs. augmented baseline & 0.902 & 0.903 & 0.001 & [-0.007, 0.010] & 0.460 & no \\
1 & Controlled & fingerprint vs. motif volume & 0.434 & 0.739 & 0.306 & [0.252, 0.351] & $<$0.001 & yes \\
1 & Prevalence & combined vs. augmented baseline & 0.773 & 0.789 & 0.016 & [-0.003, 0.036] & 0.101 & no \\
1 & Prevalence & fingerprint vs. motif volume & 0.244 & 0.526 & 0.282 & [0.220, 0.338] & $<$0.001 & yes \\
2 & Controlled & combined vs. augmented baseline & 0.902 & 0.902 & 0.001 & [-0.009, 0.010] & 0.613 & no \\
2 & Controlled & fingerprint vs. motif volume & 0.421 & 0.699 & 0.278 & [0.221, 0.325] & $<$0.001 & yes \\
2 & Prevalence & combined vs. augmented baseline & 0.722 & 0.721 & -0.001 & [-0.019, 0.017] & 0.862 & no \\
2 & Prevalence & fingerprint vs. motif volume & 0.230 & 0.464 & 0.234 & [0.159, 0.301] & $<$0.001 & yes \\
3 & Controlled & combined vs. augmented baseline & 0.750 & 0.759 & 0.008 & [-0.014, 0.031] & 0.464 & no \\
3 & Controlled & fingerprint vs. motif volume & 0.276 & 0.393 & 0.117 & [0.026, 0.200] & 0.034 & yes \\
3 & Prevalence & combined vs. augmented baseline & 0.589 & 0.615 & 0.026 & [0.000, 0.053] & 0.069 & no \\
3 & Prevalence & fingerprint vs. motif volume & 0.126 & 0.278 & 0.153 & [0.084, 0.233] & $<$0.001 & yes \\
\bottomrule
\end{tabular}
\end{table}
}{\ResPendingTable{RQ1 representation diagnostics.}{tab:res_rq1_representations}}

Table~\ref{tab:res_rq1_triangles} examines the eight triangle roles separately.
Adding them directly to the baseline changes AP by 0.0012 in the controlled
cohort and by 0.0001 in the prevalence cohort. Using all 40 roles instead of the
32 star and two-node roles changes AP by 0.0071 and 0.0041, respectively. None
of these differences is significant.

\IfFileExists{generated/results_rq1_triangles.tex}{\begin{table}[ht]
\centering
\caption{RQ1 triangle-role comparison in the 40-role Hugging Face analysis at the selected 30-day window using the default timestamp ordering. Corrected $p$-values account for the representation comparisons.}
\label{tab:res_rq1_triangles}
\footnotesize
\begin{tabular}{@{}l l r r r c r l@{}}
\toprule
\textbf{Cohort} & \textbf{Comparison} & \textbf{Ref.} & \textbf{Alt.} & \textbf{$\Delta$AP} & \textbf{95\,\% CI} & \textbf{$p_{\mathrm{BH}}$} & \textbf{Significant?} \\
\midrule
Controlled & triangles added to baseline & 0.735 & 0.736 & 0.001 & [-0.020, 0.024] & 0.540 & no \\
Controlled & 40 roles versus 32 roles & 0.752 & 0.759 & 0.007 & [-0.011, 0.026] & 0.464 & no \\
Prevalence & triangles added to baseline & 0.599 & 0.599 & 0.000 & [-0.025, 0.025] & 0.547 & no \\
Prevalence & 40 roles versus 32 roles & 0.611 & 0.615 & 0.004 & [-0.016, 0.024] & 0.448 & no \\
\bottomrule
\end{tabular}
\end{table}
}{\ResPendingTable{RQ1 triangle-role ablation.}{tab:res_rq1_triangles}}

In the binary Real-CATS comparison, the selected window is 30 days.
Baseline AP is 0.9344 and combined AP is 0.9356.
The difference is 0.0011 with a confidence interval of $[-0.0018,0.0042]$ and
$p=0.2440$. F1 decreases from 0.8481 to 0.8351. The positive AP direction is not
preserved under every timestamp ordering. The motif-only model still reaches an
AP of 0.8708.

\IfFileExists{generated/results_realcats_binary.tex}{\begin{table}[ht]
\centering
\caption{Binary comparison on the activity-matched Real-CATS sample using the default timestamp ordering.}
\label{tab:res_realcats_binary}
\footnotesize
\begin{tabular}{@{}l r@{}}
\toprule
\textbf{Quantity} & \textbf{Value} \\
\midrule
Addresses & 6,134 \\
Selected window & 30 days \\
One-standard-error window & 3 days \\
Baseline AP & 0.934 \\
Motif-only AP & 0.871 \\
Combined AP & 0.936 \\
$\Delta$AP & 0.001 \\
95\,\% CI & [-0.002, 0.004] \\
Permutation $p$-value & 0.244 \\
\bottomrule
\end{tabular}
\end{table}
}{\ResPendingTable{Binary Real-CATS comparison.}{tab:res_realcats_binary}}

For the XGBoost comparison, the AP difference is 0.0217 in the controlled cohort
($p=0.1224$) and $-0.0119$ in the prevalence cohort ($p=0.7470$).
% END PROFESSOR FEEDBACK PF19

